\chapter*{General Introduction}
\addcontentsline{toc}{chapter}{General Introduction}
\markboth{General Introduction}{General Introduction}
\label{chap:introduction}


This rapid progress in NLP over the past several years can be mostly attributed to the Transformer architecture which is capable of learning arbitrarily long and complex dependencies between words in a parallelized fashion, and has already achieved state of the art accuracy for many possible sequence modeling tasks. The problem is that this expressive capability does not come free, as the high memory requirements and immense computational throughput and energy consumption associated with training these large models effectively raise the performance ceiling to a near impossible level for practical deployment.

The issue of closing the gap between the scale used for training and that for deployment has led to the proliferation of neural network pruning as a top compression technique. The ultimate goal of pruning is to detect and eliminate either redundant parameters or structures from the already trained or during the training of network and not alter its current performance much. Pruning is, in fact, not a single problem but a highly complex optimization. One has to consider not only the architecture of the pruning, but also the condition for pruning (Whether to rely on weights magnitudes, on gradients or loss curvature), and the condition needed to make recovery.

The growing literature on pruning Transformers has revealed a somewhat noticeable problem with current state-of-the-art which is that almost all existing techniques perform pruning via a local score. By this, we mean that the importance of an individual parameter or a group of them is computed independent of any other structural components, and without the ability to know the global effect of all pruning. While it works exceptionally well in low-compression situations, performance deteriorates dramatically when we desire high sparsity. At very aggressive sparsity levels, the correct distribution of the pruning budget over the model is far more important than the precision of the local scores.

The objective of this report is to review comprehensively the state-of-the-art in Transformer pruning, analyze the existing methodologies, and formally isolate this ``allocation problem'' to motivate future search-based pruning frameworks.

To achieve this, the report is divided into two main parts. The first part establishes the foundational knowledge required for this study and spans two chapters. Chapter \ref{chap:foundations} covers the theoretical foundations of neural networks. Chapter \ref{chap:specialized_architectures} then examines the architecture of the Transformer, defines its structural components and the specific hardware-aware efficiency metrics that practically drive the need for compression.

This second part of this report covers some of the current state-of-the-art literature of Transformer pruning. Chapter \ref{chap:neural_network_pruning} introduces the formal definition of pruning in different structural regimes. In Chapter \ref{chap:related_work_pruning_unstructured} we analyze the existing work about Unstructured and Heterogeneous Search-based techniques for pruning. We mainly focus on methods that are trying to optimize flexibility and the space search at weight and layer profile level. In Chapter \ref{chap:related_work_pruning_structured} we address the problem of Structured Methods where entire units are removed at block-level, so that the result is a hardware-oblivious dense-shape. Then, Chapter \ref{chap:pruning_comparison} gives a literature overview of the comparison of existing methods between different regimes and formalizes the main shortcomings of the current evaluation procedures by clearly pointing out the limitations of local scoring algorithms in finding the solution for the global allocation problem.